% from ega4-16-fr.tex
\oldpage[IV-4]{56}
que $\Omega_{X/S}^1$ est de rang $1$, et que le morphisme $X\to S$ est
différentiellement lisse jusqu'à l'ordre $p-1$ (16.11.3, (iii)), mais non
jusqu'à l'ordre $p$. Cependant, la démonstration de (16.12.2) prouve que si
$\Omega_{X/S}^1$ est localement libre, et si $n!1_{\mathcal{O}_X}$ est
inversible dans $\Gamma(X,\mathcal{O}_X)$, alors $X$ est différentiellement
lisse sur $S$ jusqu'à l'ordre $n$.
\end{env}

% (continued in ega4-17-fr.tex, before its first page marker)
\section[Morphismes lisses, morphismes non ramifiés (ou nets) morphismes étales]{Morphismes lisses, morphismes non ramifiés (ou nets)\\morphismes étales}
\label{section:IV.17}

Dans le présent paragraphe, nous reprenons les notions étudiées dans (0, 19), exprimées à l'aide du langage géométrique des schémas et du point de vue global, pour les préschémas localement de présentation finie sur un préschéma de base donné. La plupart des résultats (à l'exception des n\textsuperscript{os} 17.7, 17.8, 17.9, 17.13 et 17.16) se réduisent d'ailleurs à des variantes de propriétés déjà rencontrées dans (0, 19). Pour des résultats plus spéciaux sur les morphismes étales, le lecteur consultera le § 18.

\subsection{Morphismes formellement lisses, morphismes formellement non ramifiés, morphismes formellement étales}
\EGAReaderTarget{ega4.c04.subsec.17-1}{subsection:IV.17.1}

\begin{definition}[17.1.1]
\EGAReaderTarget{ega4.c04.def.17-1-1}{IV.17.1.1}
Soit $f:X\to Y$ un morphisme de préschémas. On dit que $f$ est \emph{formellement lisse} (resp. \emph{formellement non ramifié} ou \emph{formellement net}, resp. \emph{formellement étale}) si, pour tout schéma affine $Y'$, tout sous-schéma fermé $Y'_0$ de $Y'$ défini par un Idéal nilpotent $\sh{J}$ de $\sh{O}_{Y'}$, et tout morphisme $Y'\to Y$, l'application
\[
\label{IV.17.1.1.1}
  \Hom_Y(Y',X)\longrightarrow\Hom_Y(Y'_0,X)
  \tag{17.1.1.1}
\]
déduite de l'injection canonique $Y'_0\to Y'$, est \emph{surjective} (resp. \emph{injective}, resp. \emph{bijective}).
\end{definition}

On dit encore alors que $X$ est \emph{formellement lisse} (resp. \emph{formellement non ramifié} ou \emph{formellement net}, resp. \emph{formellement étale}) sur $Y$.

Il est clair que dire que $f$ est formellement étale signifie qu'il est \emph{à la fois} formellement lisse et formellement non ramifié.

\begin{remarks}[17.1.2]
\label{IV.17.1.2}
\begin{enumerate}
  \item[(i)] Supposons que $Y=\Spec(A)$ et $X=\Spec(B)$ soient affines, de sorte que $f$ provient d'un homomorphisme d'anneaux $\varphi:A\to B$. En vertu de (0, 19.3.1 et 0, 19.10.1), dire que $f$ est formellement lisse (resp. formellement non ramifié, resp. formellement étale) signifie que $\varphi$ fait de $B$ une $A$-algèbre \emph{formellement lisse} (resp. \emph{formellement non ramifiée}, resp. \emph{formellement étale}) pour les topologies \emph{discrètes} sur $A$ et $B$.
  \item[(ii)] Pour vérifier que $f$ est formellement lisse (resp. formellement non ramifié, resp. formellement étale), on peut, dans la définition (17.1.1) se borner au cas où $\sh{J}^2=0$. En effet, si $f$ vérifie la condition correspondante de la définition (17.1.1) dans ce cas particulier, et si l'on a $\sh{J}^n=0$, on considère le sous-schéma fermé $Y'_j$ de $Y'$ défini par l'Idéal $\sh{J}^{j+1}$ pour $0\leq j\leq n-1$, de sorte que $Y'_j$ est un sous-schéma fermé de $Y'_{j+1}$ défini par un Idéal de carré nul; l'hypothèse entraîne que chacune des applications
  \[
    \Hom_Y(Y'_{j+1},X)\longrightarrow\Hom_Y(Y'_j,X)qquad(0\leq j\leq n-1)
  \]
  